\chapter{Presentación}

En una sociedad caracterizada por el constante desarrollo tecnológico y la transformación digital, la capacidad para resolver problemas constituye una competencia fundamental en la formación académica y profesional de las personas. La informática, más allá del uso de dispositivos tecnológicos, proporciona herramientas conceptuales y metodológicas que permiten analizar situaciones, organizar información, diseñar estrategias y construir soluciones eficientes.

Aprender a programar no consiste únicamente en memorizar instrucciones o dominar la sintaxis de un lenguaje de programación. Implica desarrollar habilidades de razonamiento lógico, pensamiento crítico, creatividad, capacidad de abstracción y organización sistemática de ideas.

En este sentido, \emph{\TituloLibro. \SubtituloLibro} surge como una propuesta educativa orientada a proporcionar al lector los conocimientos teóricos y prácticos necesarios para abordar problemas mediante soluciones computacionales.

La obra toma como fundamento directo la lógica de programación desarrollada por \citet{farrell2013}. Como bibliografía complementaria del proyecto se proponen los enfoques de resolución de problemas de Sprankle y Hubbard \citeyearpar{sprankle2012}, las estrategias de \citet{spraul2012}, el pensamiento computacional de \citet{beecher2017} y el aprendizaje mediante problemas de \citet{zingaro2021}, cuyas ediciones se identifican en las referencias.

Los ejemplos contextualizados incluyen conexiones breves con aplicaciones de inteligencia artificial, visión artificial y aprendizaje profundo.

El contenido se organiza de manera gradual, comenzando con el análisis de situaciones cotidianas y avanzando hacia el diseño de algoritmos, la construcción de programas, la realización de pruebas y la mejora de soluciones computacionales.

Se busca que el lector asuma un papel activo en su proceso de aprendizaje, mediante la exploración de problemas, la formulación de hipótesis, la experimentación y la evaluación de sus propias soluciones.
